\chapter{Introduction}

\section{Context}

Financial prices do not change at regular intervals. At short time scales, one
observes a sequence of events: an upward change, a downward change, a trade, the
addition or removal of an order. Activity alternates between calm periods and bursts.
This temporal dependence is called \emph{clustering} and is linked to volatility
\citep{bouchaud2018}.

A homogeneous Poisson process assumes a constant rate and independent increments.
This assumption is useful as a starting point, but it cannot represent the fact that
a recent change tends to be followed by further changes. The Hawkes process
\citep{hawkes1971,hawkes1971b} removes this limitation by making the intensity depend
on its own history, and it has become a reference model in high-frequency finance
\citep{bacry2015,hawkes2018}.

The practical challenge is estimating the parameters. MLE has a clear statistical
foundation, but it requires iterative optimization for each new window.
\citet{lee2023} proposes paying a training cost once and using an LSTM to map the
observed sequence directly to the parameters. This strategy is \emph{amortized
inference}: many future inferences reuse the learned knowledge \citep{cranmer2020}.

\section{Problem}

The reference paper uses United States equities between 2018 and 2019. Its NBBO data
were not published, and important details of the protocol are missing. Moreover, only
one compact recurrent architecture was studied. Crypto-assets offer a natural test
case: continuous trading, 24 hours a day, and public best bid/offer data
\citep{makarov2020}. The present work asks:

\begin{quote}
Can an attention-based architecture estimate the parameters of a 2D Hawkes process on
cryptocurrency price movements with accuracy and stability equal to or better than
Lee's LSTM, while preserving faster inference than MLE?
\end{quote}

\section{Hypotheses}

The hypotheses were registered before the final run and are kept as originally
written:

\begin{description}
  \item[H1.] Binance's public USD-M data make it possible to build sequences
  structurally equivalent to Lee's inputs.
  \item[H2.] The LSTM achieves error of the same order as MLE with substantially lower
  latency.
  \item[H3.] Attention improves error or robustness on sequences with long-range
  dependencies.
\end{description}

H3 may be rejected. A negative result, if obtained under a frozen protocol, is
scientifically informative.

\section{Objectives}

The general objective is to reproduce and extend neural Hawkes estimation on crypto.
Specific objectives:
\begin{enumerate}
  \item build and audit the top-of-book dataset;
  \item implement the 2D Hawkes process, simulation and MLE;
  \item reproduce the paper's LSTM;
  \item adapt the pipeline to BTCUSDT and ETHUSDT;
  \item design a comparable Transformer;
  \item evaluate accuracy, validity, cost and robustness;
  \item apply the best estimator to Hawkes volatility.
\end{enumerate}
Objectives 4 and 7 are addressed by the empirical stage
(Sections~\ref{sec:real-diagnostics} to~\ref{sec:real-robustness}).

\section{Contributions}

The first contribution concerns reproduction: open-source code, a record of the
paper's gaps, and a quantitative reproduction of the sampling distribution at the
published fixed point. The second is methodological: an explicit separation between
synthetic reproduction, empirical adaptation and extension, with paired statistical
tests and robust metrics that show why the aggregate MSE favors the networks. The
third is architectural: a fair comparison of LSTM and Transformer, with a negative
result for attention. The fourth is empirical: a characterization of the mid-price
events of BTCUSDT and ETHUSDT, evidence that their estimated parameters lie outside
the typical region for equities, a goodness-of-fit diagnosis that separates the effect of
the 100~ms grid from misspecification, a pre-registered evaluation of networks trained on
MLE labels of real windows, and a comparison of Hawkes volatility with realized volatility.
